\uuid{B746}
\exo7id{1643}
\titre{exo7 1643}
\auteur{barraud}
\organisation{exo7}
\datecreate{2003-09-01}
\isIndication{false}
\isCorrection{false}
\chapitre{Réduction d'endomorphisme, polynôme annulateur}
\sousChapitre{Diagonalisation}
\cpge{mpsi-mp2i.matrices-applications-lineaires.a, pcsi.matrices-applications-lineaires.a, ptsi.matrices-applications-lineaires.a, mpsi-mp2i.applications-lineaires.c, bcpst1.applications-lineaires-matrices, mp-mpi.reduction-endomorphismes.d, pc.reduction.c, psi.reduction-endomorphismes.c, pt.reduction-endomorphismes.c, mp-mpi.reduction-endomorphismes.c, pc.reduction.b, psi.reduction-endomorphismes.b, pt.reduction-endomorphismes.b, bcpst2.valeurs-propres-vecteurs-propres.b, bcpst2.valeurs-propres-vecteurs-propres.a}
\module{Algèbre}
\niveau{L2}
\difficulte{}

\contenu{
\texte{
Dans l'espace vectoriel $\R_{3}[X]$, on considère l'application linéaire
suivante~:
$$
  u~:
  \begin{array}{rcl}
    \R_{3}[X] & \rightarrow  & \R_{3}[X] \\[2ex]
    P         & \mapsto & P (0)X^{3}
                        +P'(0)X^{2}
                        +\frac{1}{2}P''(0)X
                        +\frac{1}{6}P'''(0)
  \end{array}
$$
}
\begin{enumerate}
    \item \question{Ecrire la matrice $A$ de $u$ dans la base canonique. Calculer  $A^{2}$.}
    \item \question{$u$ est-elle diagonalisable~? Si oui, donner une base de $\R^{3}[X]$
  formée de vecteurs propres de $u$.}
\end{enumerate}
}
